\section{损失函数：量化分类器的好坏}

有了线性分类器的结构 $f(x, W) = Wx + b$，下一个关键问题是：如何找到最优的 $W$？为此需要定义一个\textbf{损失函数}（loss function），也称\textbf{目标函数}（objective function），来量化当前参数 $W$ 下分类器的表现有多差。

\begin{figure}[H]
\centering
\includegraphics[width=0.95\textwidth]{slides-images/slide-065.jpg}
\caption{损失函数定义：给定训练数据集 $\{(x_i, y_i)\}_{i=1}^N$，损失函数 $L$ 衡量预测分数与真实标签之间的差距\protect\footnotemark}
\end{figure}
\footnotetext{Slides 第68页。}

\subsection{损失函数的一般形式}

给定训练数据集 $\{(x_i, y_i)\}_{i=1}^N$，其中 $x_i$ 是图像，$y_i$ 是整数标签，总损失定义为所有样本损失的平均值：

$$L = \frac{1}{N} \sum_{i=1}^{N} L_i(f(x_i, W), y_i)$$

其中 $L_i$ 是单个样本的损失函数，$f(x_i, W)$ 是模型对样本 $x_i$ 的预测分数向量。不同的 $L_i$ 定义方式产生不同的分类器。

\subsection{多类 SVM 损失（Hinge Loss）}

\begin{figure}[H]
\centering
\includegraphics[width=0.95\textwidth]{slides-images/slide-085.jpg}
\caption{多类 SVM 损失（Hinge Loss）的定义与几何解释\protect\footnotemark}
\end{figure}
\footnotetext{Slides 第88页。}

SVM 损失（也称 Hinge Loss）的核心思想是：对于每个样本 $i$，正确类别的分数 $s_{y_i}$ 应该比任何错误类别的分数 $s_j$ 至少高出一个安全边距（margin）$\Delta$（通常取 1）。

$$L_i = \sum_{j \neq y_i} \max(0, s_j - s_{y_i} + \Delta)$$

其中：
\begin{itemize}
    \item $s = f(x_i, W)$：输入 $x_i$ 的分数向量
    \item $s_{y_i}$：正确类别 $y_i$ 的分数
    \item $s_j$：其他类别 $j$ 的分数
    \item $\Delta$：安全边距，通常取 1
\end{itemize}

\begin{figure}[H]
\centering
\includegraphics[width=0.95\textwidth]{slides-images/slide-090.jpg}
\caption{Hinge Loss 的图形化解释：当正确类分数 $s_{y_i}$ 比错误类分数 $s_j$ 高出至少 1 时，损失为 0；否则损失线性增长\protect\footnotemark}
\end{figure}
\footnotetext{Slides 第93页。}

\subsubsection{SVM 损失计算示例}

\begin{figure}[H]
\centering
\includegraphics[width=0.95\textwidth]{slides-images/slide-090.jpg}
\caption{SVM 损失计算示例：三个训练样本、三个类别\protect\footnotemark}
\end{figure}
\footnotetext{Slides 第93页。}

考虑三个类别（cat, car, frog）和三个样本的分数：

\textbf{样本 1}（cat，$y_i = 0$）：分数 $s = [3.2, 5.1, -1.7]$

$$L_1 = \max(0, 5.1 - 3.2 + 1) + \max(0, -1.7 - 3.2 + 1) = \max(0, 2.9) + \max(0, -3.9) = 2.9 + 0 = 2.9$$

\textbf{样本 2}（car，$y_i = 1$）：分数 $s = [1.3, 4.9, 2.0]$

$$L_2 = \max(0, 1.3 - 4.9 + 1) + \max(0, 2.0 - 4.9 + 1) = \max(0, -2.6) + \max(0, -1.9) = 0 + 0 = 0$$

\textbf{样本 3}（frog，$y_i = 2$）：分数 $s = [2.2, 2.5, -3.1]$

$$L_3 = \max(0, 2.2 - (-3.1) + 1) + \max(0, 2.5 - (-3.1) + 1) = 6.3 + 6.6 = 12.9$$

总损失：$L = \frac{1}{3}(2.9 + 0 + 12.9) = 5.27$

\begin{knowledgebox}{关于 SVM 损失的几个重要问题}
\begin{itemize}
    \item \textbf{$L_i$ 的最小值是多少？} 答：$0$。当正确类别的分数比所有错误类别至少高出边距 $\Delta$ 时。
    \item \textbf{$L_i$ 的最大值是多少？} 答：$+\infty$。当正确类别的分数远低于错误类别时。
    \item \textbf{初始化时 $W$ 很小，所有分数接近 0，此时 $L_i \approx$ ?} 答：$C - 1$（类别数减1）。因为每个错误类别贡献约 $\max(0, 0 - 0 + 1) = 1$，共 $C-1$ 个错误类别。这是一个有用的\textbf{sanity check}。
\end{itemize}
\end{knowledgebox}

\begin{figure}[H]
\centering
\includegraphics[width=0.95\textwidth]{slides-images/slide-095.jpg}
\caption{关于 SVM 损失的思考题：如果将 $\max(0, \cdot)$ 替换为 $\max(0, \cdot)^2$，会得到不同的分类器吗？\protect\footnotemark}
\end{figure}
\footnotetext{Slides 第98页。}

将 Hinge Loss 的线性部分替换为平方：$L_i = \sum_{j \neq y_i} \max(0, s_j - s_{y_i} + 1)^2$，会得到一个不同的分类器。平方版本对大的违规（margin violation）惩罚更严厉，对小的违规惩罚更轻。选择哪种形式取决于你对错误的``容忍度''。

\subsection{Softmax 分类器（交叉熵损失）}

SVM 损失只关心正确类分数是否比错误类分数高出安全边距，不对分数赋予概率意义。Softmax 分类器则将分数转换为\textbf{概率分布}。

\subsubsection{Softmax 函数}

\begin{figure}[H]
\centering
\includegraphics[width=0.95\textwidth]{slides-images/slide-070.jpg}
\caption{Softmax 函数：先将分数指数化得到正数，再归一化得到概率分布\protect\footnotemark}
\end{figure}
\footnotetext{Slides 第73页。}

给定分数向量 $s = f(x_i, W)$，Softmax 函数将其转换为概率分布：

$$P(Y = k \mid X = x_i) = \frac{e^{s_k}}{\sum_j e^{s_j}}$$

\begin{importantbox}{Softmax 函数的两步操作}
\begin{enumerate}
    \item \textbf{指数化}：$e^{s_k}$ 确保所有值为正
    \item \textbf{归一化}：除以所有指数值之和，确保概率总和为 1
\end{enumerate}
输出可以解释为：``根据当前参数 $W$，模型认为图像 $x_i$ 属于类别 $k$ 的概率。''
\end{importantbox}

\subsubsection{交叉熵损失}

Softmax 分类器的损失函数定义为正确类别概率的\textbf{负对数}：

$$L_i = -\log P(Y = y_i \mid X = x_i) = -\log \left(\frac{e^{s_{y_i}}}{\sum_j e^{s_j}}\right)$$

\begin{figure}[H]
\centering
\includegraphics[width=0.95\textwidth]{slides-images/slide-080.jpg}
\caption{Softmax 损失定义：$L_i = -\log P(Y = y_i \mid X = x_i)$，即正确类别概率的负对数\protect\footnotemark}
\end{figure}
\footnotetext{Slides 第83页。}

\textbf{直觉理解}：
\begin{itemize}
    \item 我们希望\textbf{最大化}正确类别的概率 $P(Y = y_i | X = x_i)$
    \item 等价于\textbf{最小化}其负对数 $-\log P$
    \item 取负号将最大化问题转为最小化问题
    \item 取对数使数值更易处理（对数函数是单调递增的，不改变最优解）
\end{itemize}

\subsubsection{交叉熵损失计算示例}

以之前的例子（cat，分数 $s = [3.2, 5.1, -1.7]$）：

\textbf{Step 1：指数化}
$$e^{3.2} = 24.5, \quad e^{5.1} = 164.0, \quad e^{-1.7} = 0.18$$

\textbf{Step 2：归一化}
$$P(\text{cat}) = \frac{24.5}{24.5 + 164.0 + 0.18} = \frac{24.5}{188.68} \approx 0.13$$

\textbf{Step 3：计算损失}
$$L_i = -\log(0.13) \approx 2.04$$

当前模型认为这张猫的图片是猫的概率只有 13\%——说明 $W$ 还需要优化。

\begin{knowledgebox}{交叉熵损失的取值范围与 Sanity Check}
\begin{itemize}
    \item \textbf{最小值}：$0$（当正确类别的概率为 1 时，$-\log(1) = 0$）
    \item \textbf{最大值}：$+\infty$（当正确类别的概率趋近于 0 时）
    \item \textbf{Sanity Check}：初始化时所有分数接近相等，每个类别概率约为 $1/C$，此时 $L_i \approx -\log(1/C) = \log(C)$。对于 CIFAR-10（$C=10$），$L_i \approx \ln(10) \approx 2.3$。如果初始损失远大于这个值，说明实现可能有 bug。
\end{itemize}
\end{knowledgebox}

\subsubsection{交叉熵的多种等价解释}

\begin{figure}[H]
\centering
\includegraphics[width=0.85\textwidth]{slides-images/slide-078.jpg}
\caption{交叉熵损失的多种解释：最大似然估计、KL 散度、交叉熵\protect\footnotemark}
\end{figure}
\footnotetext{Slides 第81页。}

同一个损失函数有多种等价的理论解释：

\begin{enumerate}
    \item \textbf{最大似然估计（MLE）}：选择使观测数据出现概率最大的参数 $W$
    \item \textbf{KL 散度最小化}：最小化真实分布 $p$（one-hot）与预测分布 $q$（softmax 输出）之间的 KL 散度
    \item \textbf{交叉熵最小化}：$H(p, q) = H(p) + D_{KL}(p \| q)$。由于 one-hot 分布的熵 $H(p) = 0$，所以 $H(p,q) = D_{KL}(p \| q) = -\log q_{y_i}$
\end{enumerate}

\begin{knowledgebox}{Softmax 分类器 = 多项逻辑回归}
Softmax 分类器本质上就是\textbf{多项逻辑回归}（Multinomial Logistic Regression）。如果你在其他课程中学过二分类的逻辑回归，Softmax 是它在多分类场景下的自然推广。在深度学习框架中，这个损失函数通常被称为 \texttt{CrossEntropyLoss} 或 \texttt{BCE}（Binary Cross Entropy，用于二分类）。
\end{knowledgebox}

\subsection{SVM 损失 vs Softmax 损失}

\begin{figure}[H]
\centering
\includegraphics[width=0.95\textwidth]{slides-images/slide-100.jpg}
\caption{Softmax vs SVM 损失对比：给定相同分数，两种损失的行为差异\protect\footnotemark}
\end{figure}
\footnotetext{Slides 第103页。}

\begin{table}[H]
\centering
\begin{tabular}{lcc}
\toprule
\textbf{特性} & \textbf{SVM 损失} & \textbf{Softmax 损失} \\
\midrule
公式 & $\sum_{j \neq y_i} \max(0, s_j - s_{y_i} + 1)$ & $-\log\left(\frac{e^{s_{y_i}}}{\sum_j e^{s_j}}\right)$ \\
分数解释 & 无特殊含义 & 概率（经 softmax 归一化后） \\
关注点 & 正确类是否比错误类高出边距 & 正确类的概率有多高 \\
满足条件后 & 损失为 0，不再优化 & 永远不会精确为 0，持续优化 \\
最小值 & 0 & 0（理论极限） \\
初始化 sanity check & $C - 1$ & $\log(C)$ \\
\bottomrule
\end{tabular}
\caption{SVM 损失与 Softmax 损失的对比}
\end{table}

\begin{figure}[H]
\centering
\includegraphics[width=0.95\textwidth]{slides-images/slide-101.jpg}
\caption{思考题：当正确类分数从 10 增加到 20 时，SVM 损失和 Softmax 损失的变化\protect\footnotemark}
\end{figure}
\footnotetext{Slides 第104页。}

\begin{importantbox}{两种损失的核心差异}
\textbf{SVM 损失}只关心``正确类是否够好''——一旦正确类分数超过所有错误类分数加上边距，损失就为 0，不再有进一步优化的动力。\\
\textbf{Softmax 损失}永远在优化——即使正确类概率已经很高（如 99\%），损失仍然大于 0，模型会继续尝试推高正确类的概率。
\end{importantbox}

\subsection{本章小结}

损失函数是连接``分类器结构''和``参数优化''的桥梁。本节介绍了两种重要的损失函数：
\begin{itemize}
    \item \textbf{SVM/Hinge 损失}：基于边距的思想，要求正确类分数比错误类至少高出 $\Delta$
    \item \textbf{Softmax/交叉熵损失}：将分数转化为概率，最大化正确类别的概率（等价于最小化 KL 散度/交叉熵）
\end{itemize}

两种损失在实践中都被广泛使用，Softmax 损失在现代深度学习中更为常见。下一讲将介绍如何通过\textbf{优化算法}（如梯度下降）来找到使损失函数最小的参数 $W$。

%% ========== 总结与延伸 ==========

